\begin{definition}
    Define the \textbf{exponent} $\forall a \in \bZ, b \in \bZ \cup \{0\}$ as \[
    a^b = \begin{cases}
        1 & \text{if } b=0\\
        a \cdot a^{b-1} & \text{otherwise}
    \end{cases}
    \]
\end{definition}
For properties of exponents, see Appendix B.

\begin{definition}[The polynomial ring $\Qx$]
We define a polynomial $f(x) \in \Qx$ with degree $n$ if it can be written as
\[
f(x)=\sum_{0}^{n} a_i x^i
\]
where $n\in\mathbb{Z}_{\geq0}$ and each coefficient $a_i$ lies in $\mathbb{Q}$ and $a_n \neq 0$.
\end{definition}

\subsection{Ring Axioms}

\begin{proof} %%needs to be fixed is kinda wishy washy
We check the ring conditions.  Addition stays inside $\Qx$ because each
coefficient $a_i+b_i$ is rational.  Multiplication also stays inside $\Qx$
because each coefficient $c_k$ is built from finitely many sums and products of
rational numbers.

The zero polynomial is the additive identity because adding zero to every
coefficient changes nothing.  The polynomial $-f(x)$ is the additive inverse of
$f(x)$ because each coefficient becomes $a_i+(-a_i)=0$.  The constant polynomial
$1$ is the multiplicative identity because multiplying by $1$ leaves the
coefficients unchanged.

The associative and commutative laws come from the same laws in $\mathbb{Q}$,
coefficient by coefficient.  For instance, addition is associative because
\[
(a_k+b_k)+c_k=a_k+(b_k+c_k)
\]
for each coefficient.  The multiplication laws and distributivity are checked in
the same way from the rational-number laws.

Finally, the zero polynomial has constant coefficient $0$, while the constant
polynomial $1$ has constant coefficient $1$.  Since $0\neq1$ in $\mathbb{Q}$,
have $0_{\Qx}\neq1_{\Qx}$.
\end{proof}

\begin{definition}
    For $f(x) \in \Qx$, let the \textbf{norm}, or $\| f(x)\|$ be equal to \[
    \| f(x) \|  = 2^{\deg{f(x)}}
    \]
with the following properties:
\begin{itemize}
    \item $\| f(x)\| \cdot\| g(x)\| = \| f(x)g(x)\|$ 
    \begin{proof}
        By Multiplied Sums, if $f(x)=\sum_{i=0}^{n} a_i x^i$ and $g(x)=\sum_{j=0}^{m} b_j x^j$, with $n$ and $m$ the degrees of $f$ and $g$ respectively, then 
        \begin{align*}
        f(x)g(x)&=\sum_{i=0}^{n} \sum_{j=0}^{m} a_i b_j x^i x^j=\sum_{i=0}^{n} \sum_{j=0}^{m} a_i b_j x^{i+j}
        \end{align*}
        Notice that to find the degree of $f(x)g(x)$, we take the maximum value of $i+j$ which is their respective maximums, $n+m$, so:
        \[
        \|f(x)g(x)\|=2^{n+m}=2^n2^m = \|f(x)g(x)\|
        \]
    \end{proof}
    \item For any $c(x) \in \Qx ^ {\times}$, $\|c(x)\| = 1$.
\begin{proof} Since $\deg{c(x)} = 0$, $\|c(x)\|=2^0 = 1$.\end{proof}
\end{itemize}
\end{definition}

% \begin{lemma}[Multiplication] i commented this out bc i kinda did it with set notation up there

% %%%TODO: this is not rigorous (use set not)
% For multiplication, if
% \[
% f(x)=a_0+a_1x+\cdots+a_mx^m
% \qquad\text{and}\qquad
% g(x)=b_0+b_1x+\cdots+b_nx^n,
% \]
% then
% \[
% f(x)g(x)=c_0+c_1x+\cdots+c_{m+n}x^{m+n},
% \]
% where each coefficient $c_k$ is found by adding all products $a_ib_j$ for which
% the exponents satisfy $i+j=k$.  For example,
% \[
% c_0=a_0b_0,\qquad c_1=a_0b_1+a_1b_0,\qquad
% c_{m+n}=a_mb_n.
% \]
% \end{lemma}

\begin{definition}[Zero, one, and negatives]
The zero element of $\Qx$ is the polynomial $0_{Qx} = 0$. Define the degree of $0_{Qx}$ to be $-\infty$.  
The one element of $\Qx$ is the constant polynomial
\[
1_{\Qx}=1 
\]
If
\[
f(x)=\sum_{0}^{n} a_i x^i
\]
then
\[
-f(x)=\sum_{0}^{n} -a_i x^i
\]
\end{definition}

\begin{theorem}[$\Qx$ has no zero divisors]
If $f(x),g(x)\in\Qx$ and
\[
f(x)g(x)=0,
\]
then $f(x)=0$ or $g(x)=0$.
\end{theorem}

\begin{proof}
Suppose, for contradiction, that
\[
f(x)g(x)=0
\]
with $f(x)\neq0$ and $g(x)\neq0$.

Since $g(x)\neq0$, by the norm lemma,
\[
\|f(x)\|\leq \|f(x)g(x)\|.
\]
Since $f(x)\neq0$, we have
\[
\|f(x)\|=2^{\deg(f(x))}>0.
\]
So $\|f(x)g(x)\|$ is positive. In particular, $f(x)g(x)\neq0$, since the norm
is only positive for nonzero polynomials.

This contradicts $f(x)g(x)=0$. Therefore at least one of $f(x)$ and $g(x)$
must be zero.
\end{proof}

\begin{corollary}[Cancellation]
If $h(x)\neq0$ and
\[
h(x)f(x)=h(x)g(x),
\]
then $f(x)=g(x)$.
\end{corollary}

\begin{proof}
Subtract the right side from the left:
\[
h(f-g)=0.
\]
Since $h\neq0$ and $\Qx$ has no zero divisors, we must have
$f-g=0$.  Thus $f=g$.
\end{proof}

\subsection{Basic Arithmetic}

\begin{lemma}[Uniqueness of zero and one]
The additive identity and multiplicative identity in $\Qx$ are unique.
In particular, if $f(x)+z(x)=f(x)$ for every $f(x)\in\Qx$, then
$z(x)=0$.  Similarly, if $f(x)u(x)=f(x)$ for every $f(x)\in\Qx$, then
$u(x)=1$.
\end{lemma}

\begin{proof}
Suppose $z$ acts like an additive identity for $f$, so $f+z=f$.  Add $-f$ to
both sides:
\[
(-f)+(f+z)=(-f)+f.
\]
By associativity and the definition of additive inverse, this becomes
\[
0+z=0.
\]
Since $0$ is the additive identity, $z=0$.

Now suppose $u$ acts like a multiplicative identity for every polynomial.  Test
this on the constant polynomial $1$:
\[
1\cdot u=1.
\]
But $1\cdot u=u$.  Hence $u=1$.
\end{proof}

\begin{lemma}[Uniqueness of additive inverses]
If $u(x)$ and $v(x)$ are both additive inverses of $f(x)$, then $u(x)=v(x)$.
\end{lemma}

\begin{proof}
Assume
\[
f+u=0
\qquad\text{and}\qquad
f+v=0.
\]
Then $f+u=f+v$.  Add $-f$ to both sides:
\[
(-f)+(f+u)=(-f)+(f+v).
\]
By associativity,
\[
((-f)+f)+u=((-f)+f)+v.
\]
Thus
\[
0+u=0+v,
\]
so $u=v$.
\end{proof}

\begin{lemma}[Zero and negatives]
For all $f(x),g(x)\in\Qx$,
\[
f(x)\cdot0=0,\qquad -(-f(x))=f(x),\qquad (-f(x))(-g(x))=f(x)g(x).
\]
\end{lemma}

\begin{proof}
First,
\[
f\cdot0=f(0+0)=f\cdot0+f\cdot0.
\]
Adding the additive inverse of $f\cdot0$ to both sides gives $f\cdot0=0$.

Next, since
\[
(-f)+f=0,
\]
the polynomial $f$ is an additive inverse of $-f$.  By uniqueness of additive
inverses,
\[
-(-f)=f.
\]

For the product of negatives, first observe that
\[
fg+(-f)g=(f+(-f))g=0\cdot g=0.
\]
Therefore $(-f)g=-(fg)$.  Applying this with $g$ replaced by $-g$ gives
\[
(-f)(-g)=-(f(-g)).
\]
Also,
\[
fg+f(-g)=f(g+(-g))=f\cdot0=0,
\]
so $f(-g)=-(fg)$.  Hence
\[
(-f)(-g)=-(-(fg))=fg.
\]
\end{proof}
%%TODO: \|f(x)\| \leq \|f(x)g(x)\| for g(x) \neq 0. (Done)
\begin{lemma}
Let $f(x),g(x)\in\Qx$ with $f(x)\neq0$ and $g(x)\neq0$. Then
\[
\|f(x)\|\leq \|f(x)g(x)\|.
\]
\end{lemma}

\begin{proof}
Let
\[
\deg(f)=m
\qquad\text{and}\qquad
\deg(g)=n.
\]
Since $g(x)\neq0$, we have $n\geq0$. Also, since $f(x)\neq0$ and
$g(x)\neq0$, the leading coefficient of $f(x)g(x)$ is the product of the
leading coefficients of $f(x)$ and $g(x)$, which is nonzero in $\mathbb{Q}$.
Thus
\[
\deg(fg)=\deg(f)+\deg(g)=m+n.
\]
Therefore
\[
\|f(x)g(x)\|=2^{\deg(fg)}=2^{m+n}.
\]
Since $n\geq0$,
\[
2^m\leq2^{m+n}.
\]
So
\[
\|f(x)\|=2^m\leq2^{m+n}=\|f(x)g(x)\|.
\]
\end{proof}


\subsection{Divisibility}

\begin{definition}[Divisibility]
For $f(x),g(x)\in\Qx$, we say $f(x)$ divides $g(x)$, written
\[
f(x)\mid g(x),
\]
if there exists $q(x)\in\Qx$ such that
\[
g(x)=f(x)q(x).
\]
\end{definition}

\begin{definition}[Greatest Common Divisor]
    For $f(x),g(x)\in\Qx$, we defince $\gcd(f(x),g(x)) = d(x)$ to satisfy:
    \begin{itemize}
        \item $d(x) \mid f(x)$ and $d(x) \mid g(x)$
        \item $\forall h(x) \in \Qx$ s.t. $h(x) \mid f(x)$ and $h(x) \mid g(x)$, $\|h(x)\| \leq \|d(x)\|$
    \end{itemize}
\end{definition} %%todo: existence and uniqueness (DONE)

\begin{theorem}[Existence and uniqueness of gcd]
Let $f(x),g(x)\in\Qx$, not both zero. Then there exists a unique monic
polynomial $d(x)\in\Qx$ such that
\[
d(x)=\gcd(f(x),g(x)).
\]
\end{theorem}

\begin{proof}
Let
\[
S=\{\,a(x)f(x)+b(x)g(x):a(x),b(x)\in\Qx,\ a(x)f(x)+b(x)g(x)\neq0\,\}.
\]
Since $f(x)$ and $g(x)$ are not both zero, the set $S$ is nonempty. Also, every
element of $S$ has positive norm. So the set
\[
N=\{\,\|s(x)\|:s(x)\in S\,\}
\]
is a nonempty subset of $\mathbb{Z}^+$. By the well-ordering principle, $N$ has
a least element. Choose $d(x)\in S$ with smallest norm. Then for some
$a(x),b(x)\in\Qx$,
\[
d(x)=a(x)f(x)+b(x)g(x).
\]

We first show that $d(x)\mid f(x)$. By the division algorithm, there exist
$q(x),r(x)\in\Qx$ such that
\[
f(x)=q(x)d(x)+r(x),
\]
where either $r(x)=0$ or
\[
\|r(x)\|<\|d(x)\|.
\]
Rearranging gives
\[
r(x)=f(x)-q(x)d(x).
\]
Since $d(x)=a(x)f(x)+b(x)g(x)$, we get
\[
r(x)=f(x)-q(x)(a(x)f(x)+b(x)g(x)).
\]
Thus
\[
r(x)=(1-q(x)a(x))f(x)+(-q(x)b(x))g(x).
\]
So if $r(x)\neq0$, then $r(x)\in S$, but it has smaller norm than $d(x)$,
which contradicts how $d(x)$ was chosen. Therefore $r(x)=0$, and hence
\[
d(x)\mid f(x).
\]
The same argument with $g(x)$ in place of $f(x)$ shows that
\[
d(x)\mid g(x).
\]

Now let $h(x)\in\Qx$ be any common divisor of $f(x)$ and $g(x)$. Then
\[
f(x)=h(x)u(x)
\qquad\text{and}\qquad
g(x)=h(x)v(x)
\]
for some $u(x),v(x)\in\Qx$. Hence
\[
d(x)=a(x)f(x)+b(x)g(x)
=a(x)h(x)u(x)+b(x)h(x)v(x).
\]
Factoring gives
\[
d(x)=h(x)(a(x)u(x)+b(x)v(x)).
\]
Thus
\[
h(x)\mid d(x).
\]
So $d(x)$ satisfies the gcd property.

Finally, make $d(x)$ monic. Let $c\in\mathbb{Q}$ be the leading coefficient of
$d(x)$. Since $d(x)\neq0$, we have $c\neq0$. Define
\[
d_0(x)=c^{-1}d(x).
\]
Then $d_0(x)$ is monic. Since $c$ is a nonzero rational number, multiplying by
$c$ or $c^{-1}$ does not change divisibility in $\Qx$. Therefore $d_0(x)$ is
also a greatest common divisor of $f(x)$ and $g(x)$.

It remains to prove uniqueness. Suppose $d_1(x)$ and $d_2(x)$ are both monic
greatest common divisors of $f(x)$ and $g(x)$. Since $d_1(x)$ is a gcd and
$d_2(x)$ is a common divisor, we have
\[
d_2(x)\mid d_1(x).
\]
Similarly,
\[
d_1(x)\mid d_2(x).
\]
So there exist $u(x),v(x)\in\Qx$ such that
\[
d_1(x)=d_2(x)u(x)
\qquad\text{and}\qquad
d_2(x)=d_1(x)v(x).
\]
Taking degrees gives
\[
\deg(d_1)=\deg(d_2)+\deg(u)
\]
and
\[
\deg(d_2)=\deg(d_1)+\deg(v).
\]
Thus $\deg(u)=\deg(v)=0$, so $u(x)$ and $v(x)$ are nonzero constant
polynomials. Since $d_1(x)$ and $d_2(x)$ are both monic, the constant $u(x)$
must be $1$. Therefore
\[
d_1(x)=d_2(x).
\]
Hence the monic greatest common divisor exists and is unique.
\end{proof}
